\section{Introdu\c{c}\~ao}\label{sec:intro}

% Texto original do grupo, preservado. Citacoes convertidas para
% biblatex (\parencite / \textcite) e paragrafo-roteiro final ajustado
% para a nova estrutura do documento.

As estruturas de concreto armado constituem a espinha dorsal dos sistemas de infraestrutura crítica; entretanto, sua durabilidade a longo prazo é fortemente afetada por mecanismos de degradação, dentre os quais se destaca a corrosão das armaduras induzida pela carbonatação. A carbonatação reduz a alcalinidade do concreto, viabilizando o início da corrosão e a perda progressiva de capacidade resistente, o que pode comprometer a segurança, a funcionalidade e o desempenho ao longo do ciclo de vida \parencite{tuutti1982, possan2010}. À medida que as redes de infraestrutura envelhecem e os recursos de manutenção se tornam cada vez mais restritos, a capacidade de prever com acurácia a vida útil remanescente (\textit{Remaining Useful Life} -- RUL) de estruturas de concreto armado tornou-se um desafio central da engenharia estrutural e da gestão de ativos informada por risco \parencite{ellingwood2005, frangopol2011, si2011}.

As abordagens tradicionais de previsão de vida útil apoiam-se em modelos físicos de profundidade de carbonatação e em formulações determinísticas de estado limite. Embora tais modelos forneçam valiosa compreensão mecanicista, eles frequentemente não capturam os efeitos combinados da variabilidade dos materiais, da incerteza ambiental, da deterioração estocástica e do erro epistêmico de modelo \parencite{possan2010, melchers2018}. Os avanços recentes em modelagem probabilística de durabilidade têm, portanto, enfatizado a integração entre quantificação de incertezas, processos estocásticos de deterioração e calibração orientada por dados \parencite{jcss2001, fib2006}.

Em paralelo, técnicas de aprendizado de máquina (\textit{Machine Learning} -- ML) e de metamodelagem têm recebido crescente atenção para a previsão da profundidade de carbonatação, da degradação estrutural e do desempenho ao longo do ciclo de vida de sistemas de concreto armado. Estudos comparativos demonstraram a capacidade de redes neurais artificiais e de outros modelos de ML de superar formulações convencionais de regressão quando treinados com dados experimentais e de campo \parencite{couto2025}. Contudo, modelos puramente orientados por dados costumam apresentar capacidade limitada de extrapolação, baixa interpretabilidade e sensibilidade à escassez de dados de treinamento, o que restringe sua aplicação direta em avaliações de confiabilidade de sistemas críticos.

No contexto da quantificação de incertezas, metamodelos consolidados como os Processos Gaussianos (GP) \parencite{rasmussen2006} e as Expansões em Caos Polinomial (PCE) \parencite{xiu2002, blatman2011} são padrão para simuladores \textit{determinísticos}, nos quais uma única avaliação caracteriza a resposta. Esses metamodelos, entretanto, não podem ser aplicados diretamente a \textit{simuladores estocásticos}, nos quais a resposta para um mesmo vetor de entrada é uma variável aleatória cuja distribuição precisa ser caracterizada. Essa é exatamente a natureza do problema de degradação por carbonatação: para uma mesma configuração de projeto, variáveis latentes de material e de ambiente produzem uma distribuição de margens de segurança, e não um valor único. Para esse cenário, foram propostos recentemente os emuladores estocásticos, com destaque para os modelos \textit{Generalized Lambda Models} (GLaM), nos quais a distribuição condicional da resposta é aproximada pela Distribuição Lambda Generalizada (GLD) e seus quatro parâmetros são representados por PCE das variáveis de entrada \parencite{ZhuSudret2020, ZhuSudret2021}.

Apesar desses avanços, identifica-se uma lacuna na literatura: os trabalhos existentes sobre confiabilidade dependente do tempo de estruturas sujeitas à carbonatação concentram-se em emular momentos estatísticos, o índice de confiabilidade $\beta(t)$ ou a profundidade média de carbonatação, e o arcabouço GLaM original é formulado para um instante fixo, sem tratamento contínuo da variável tempo. Nenhuma abordagem disponível fornece, de forma simultânea, (i) a distribuição condicional \textit{completa} da função de estado limite $g(t)$, (ii) a interpolação contínua dessa distribuição entre instantes simulados e (iii) a tradução direta dessa informação em uma distribuição probabilística de RUL apta a subsidiar decisões de manutenção.

Para preencher essa lacuna, o presente estudo propõe um emulador estocástico dependente do tempo no qual a resposta do simulador é modelada por uma GLD cujos parâmetros $\boldsymbol{\lambda}$ são funções das variáveis de entrada, representadas por PCE, e a evolução temporal desses parâmetros é aprendida por um modelo de ML que recebe o tempo como variável de entrada contínua. Uma vez treinado, o emulador fornece instantaneamente a função densidade de probabilidade da margem de segurança $g(\mathbf{X}, t)$ em qualquer instante $t$ do horizonte de análise, da qual se derivam a probabilidade de falha, a distribuição do tempo até a falha, a RUL probabilística e métricas de risco de cauda (VaR e CVaR). A metodologia é aplicada à previsão da vida útil remanescente de elementos de concreto armado sujeitos à corrosão induzida por carbonatação, integrando o modelo de avanço da frente de carbonatação de \textcite{possan2010}, a perda de seção e de aderência das armaduras \parencite{val1997, azad2010, peng2016} e a verificação do estado limite último de flexão segundo a \textcite{nbr6118}. Para verificar o arcabouço metodológico, adota-se uma versão adaptada de um problema de referência de confiabilidade, seguindo a estratégia geral de validação proposta por \textcite{pires2024}.

As principais contribuições deste trabalho são: (i) a extensão do arcabouço GLaM para o domínio temporal, por meio de uma camada de aprendizado de máquina que interpola continuamente os parâmetros da GLD entre os instantes simulados; (ii) a integração da deterioração por carbonatação em um arcabouço probabilístico de previsão de RUL que opera sobre a distribuição completa da função de estado limite, e não apenas sobre seus momentos; e (iii) a demonstração de ganhos de eficiência computacional de ordens de grandeza em relação à abordagem convencional de Monte Carlo, viabilizando análises de RUL em tempo quase real, inclusive em cenários com intervenções de manutenção. Os resultados visam a contribuir para a avaliação de ciclo de vida baseada em confiabilidade de infraestruturas de concreto armado e para o suporte a decisões de manutenção informadas por risco.

O restante do artigo está organizado da seguinte forma. A Seção~\ref{sec:emuladores} introduz o conceito de simulador estocástico e o núcleo do emulador proposto --- a Distribuição Lambda Generalizada e seu acoplamento com a Expansão em Caos Polinomial. A Seção~\ref{sec:ml} apresenta os fundamentos de aprendizado de máquina empregados, com ênfase nas redes neurais artificiais que constituem a camada temporal do emulador. A Seção~\ref{sec:metodologia} detalha a metodologia: os modelos físicos de carbonatação e de resistência residual, a definição da função de estado limite dependente do tempo, o procedimento de construção do emulador e a ponte entre a distribuição emulada e a RUL probabilística. A Seção~\ref{sec:resultados} apresenta os exemplos numéricos, iniciando pela verificação em um problema de referência e avançando para a previsão de RUL de elementos de concreto armado sob carbonatação. A Seção~\ref{sec:conclusoes} resume as conclusões.
